\section{Operational calculus}
\label{sec:operational-calculus-mealy}

The component description of a Mealy morphism can be presented as a formal
typed operational calculus. This section records that calculus in
definition--proposition--theorem form. The purpose is not to introduce new
structure, but to give a precise judgemental presentation of the data and laws
already defined.

Fix a Mealy morphism
\[
    (P,\delta,\omega):\mathbb A\rightsquigarrow\mathbb B.
\]
We use context objects
\[
    \Gamma\in\mathcal E
\]
to make the generalized-element notation formal.

\begin{definition}[Judgements]
\label{def:mealy-operational-judgements}
    The operational calculus has the following judgement forms.

    \begin{enumerate}
        \item An input-object judgement
        \[
            \Gamma\vdash u:A_0
        \]
        is interpreted as a morphism \(u:\Gamma\to A_0\).

        \item An input-arrow judgement
        \[
            \Gamma\vdash a:u\to_{\mathbb A}v
        \]
        is interpreted as a morphism \(a:\Gamma\to A_1\) such that
        \[
            s_Aa=u,
            \qquad
            t_Aa=v.
        \]
        Diagrammatically,
\[
\begin{tikzcd}[column sep=large, row sep=large]
    &
    \Gamma
        \arrow[d, "a" description]
        \arrow[dl, "u"']
        \arrow[dr, "v"]
    &
    \\
    A_0
    &
    A_1
        \arrow[l, "s_A"]
        \arrow[r, "t_A"']
    &
    A_0 .
\end{tikzcd}
\]
        records the equations \(s_Aa=u\) and \(t_Aa=v\).

        \item An output-object judgement
        \[
            \Gamma\vdash y:B_0
        \]
        is interpreted as a morphism \(y:\Gamma\to B_0\).

        \item An output-arrow judgement
        \[
            \Gamma\vdash \beta:y'\to_{\mathbb B}y
        \]
        is interpreted as a morphism \(\beta:\Gamma\to B_1\) such that
        \[
            s_B\beta=y',
            \qquad
            t_B\beta=y.
        \]

        \item A state judgement
        \[
            \Gamma\vdash x:u\leadsto y
        \]
        is interpreted as a morphism \(x:\Gamma\to P\) such that
        \[
            px=u,
            \qquad
            qx=y.
        \]
        Thus \(x\) is a generalized state whose input interface is \(u\) and
        whose output interface is \(y\):
 \[
\begin{tikzcd}[column sep=large, row sep=large]
    &
    \Gamma
        \arrow[d, "x" description]
        \arrow[dl, "u"']
        \arrow[dr, "y"]
    &
    \\
    A_0
    &
    P
        \arrow[l, "p"]
        \arrow[r, "q"']
    &
    B_0 .
\end{tikzcd}
\]    \end{enumerate}
\end{definition}

\begin{definition}[One-step execution]
\label{def:one-step-execution}
    Suppose
    \[
        \Gamma\vdash a:u\to_{\mathbb A}v
        \qquad
        \text{and}
        \qquad
        \Gamma\vdash x:v\leadsto y.
    \]
    These judgements determine a morphism
    \[
        \langle a,x\rangle:
        \Gamma\to A_1\times_{t_A,A_0,p}P
    \]
    by the pullback square
    \[
        \begin{tikzcd}
            A_1\times_{A_0}P
                \arrow[r, "\pi_P"]
                \arrow[d, "\pi_{A_1}"']
                \arrow[dr, phantom, "\lrcorner", very near start]
            &
            P \arrow[d, "p"]
            \\
            A_1 \arrow[r, "t_A"']
            &
            A_0 .
        \end{tikzcd}
    \]
    The one-step execution judgement
    \[
        \Gamma\vdash
        x\xrightarrow{\;a/\beta\;}_{P}x'
    \]
    is defined by
    \[
        x':=\delta\langle a,x\rangle,
        \qquad
        \beta:=\omega\langle a,x\rangle.
    \]
    The intended picture is
    \[
        \begin{tikzcd}[column sep=large]
            x
                \arrow[r, "{a/\beta}"]
            &
            x' .
        \end{tikzcd}
    \]
\end{definition}

\begin{proposition}[Typing of one-step execution]
\label{prop:typing-one-step-execution}
    Under the hypotheses of Definition~\ref{def:one-step-execution}, the
    conclusion is well typed:
    \[
        \Gamma\vdash x':u\leadsto y'
    \]
    and
    \[
        \Gamma\vdash \beta:y'\to_{\mathbb B}y,
    \]
    where
    \[
        y':=qx'.
    \]
\end{proposition}

\begin{proof}
    \leavevmode
    \begin{enumerate}
        \item \textbf{Pull back the transition typing equation.}
        The Mealy typing equation
        \[
            p\delta=s_A\pi_{A_1}
        \]
        pulled back along
        \[
            \langle a,x\rangle:\Gamma\to A_1\times_{A_0}P
        \]
        gives
        \[
            px'=s_Aa=u.
        \]

        \item \textbf{Pull back the output source equation.}
        The Mealy typing equation
        \[
            s_B\omega=q\delta
        \]
        gives
        \[
            s_B\beta=qx'=y'.
        \]

        \item \textbf{Pull back the output target equation.}
        The Mealy typing equation
        \[
            t_B\omega=q\pi_P
        \]
        gives
        \[
            t_B\beta=qx=y.
        \]

        \item \textbf{Conclude well-typedness.}
        Hence
        \[
            \Gamma\vdash x':u\leadsto y',
        \]
        and
        \[
            \Gamma\vdash \beta:y'\to_{\mathbb B}y.
        \]
    \end{enumerate}
\end{proof}

\begin{definition}[Identity execution rule]
\label{def:identity-execution-rule}
    The identity execution rule is the admissible rule
    \[
        \frac{
            \Gamma\vdash x:v\leadsto y
        }
        {
            \Gamma\vdash
            x\xrightarrow{\;1_v/1_y\;}_{P}x
        }.
    \]
    Here
    \[
        1_v:=e_Av,
        \qquad
        1_y:=e_By.
    \]
\end{definition}

\begin{proposition}[Soundness of identity execution]
\label{prop:soundness-identity-execution}
    The identity execution rule is sound.
\end{proposition}

\begin{proof}
    \leavevmode
    \begin{enumerate}
        \item \textbf{Apply the Mealy unit law.}
        The Mealy unit law gives
        \[
            \delta(e_A(px),x)=x,
            \qquad
            \omega(e_A(px),x)=e_B(qx).
        \]

        \item \textbf{Translate into judgemental notation.}
        If
        \[
            \Gamma\vdash x:v\leadsto y,
        \]
        then \(px=v\) and \(qx=y\). Therefore
        \[
            \delta(e_Av,x)=x,
            \qquad
            \omega(e_Av,x)=e_By.
        \]
        This is exactly
        \[
            \Gamma\vdash
            x\xrightarrow{\;1_v/1_y\;}_{P}x.
        \]
    \end{enumerate}
\end{proof}

\begin{definition}[Sequential execution rule]
\label{def:sequential-execution-rule}
    Suppose
    \[
        \Gamma\vdash a:u\to_{\mathbb A}v,
        \qquad
        \Gamma\vdash b:v\to_{\mathbb A}w,
        \qquad
        \Gamma\vdash x:w\leadsto y.
    \]
    If
    \[
        \Gamma\vdash
        x\xrightarrow{\;b/\beta\;}_{P}x_1
    \]
    and
    \[
        \Gamma\vdash
        x_1\xrightarrow{\;a/\gamma\;}_{P}x_2,
    \]
    then the sequential execution rule derives
    \[
        \Gamma\vdash
        x\xrightarrow{\;b\circ a/\,\beta\circ\gamma\;}_{P}x_2.
    \]
    The rule is pictured as
    \[
        \begin{tikzcd}[column sep=large]
            x
                \arrow[r, "{b/\beta}"]
            &
            x_1
                \arrow[r, "{a/\gamma}"]
            &
            x_2
        \end{tikzcd}
        \qquad\leadsto\qquad
        \begin{tikzcd}[column sep=large]
            x
                \arrow[r, "{b\circ a\,/\,\beta\circ\gamma}"]
            &
            x_2 .
        \end{tikzcd}
    \]
\end{definition}

\begin{proposition}[Soundness of sequential execution]
\label{prop:soundness-sequential-execution}
    The sequential execution rule is sound.
\end{proposition}

\begin{proof}
    \leavevmode
    \begin{enumerate}
        \item \textbf{Expand the first execution.}
        By definition of one-step execution,
        \[
            x_1=\delta(b,x),
            \qquad
            \beta=\omega(b,x).
        \]

        \item \textbf{Expand the second execution.}
        Again by definition,
        \[
            x_2=\delta(a,x_1)=\delta(a,\delta(b,x)),
            \qquad
            \gamma=\omega(a,\delta(b,x)).
        \]

        \item \textbf{Use the Mealy multiplication law for states.}
        The state multiplication law gives
        \[
            \delta(b\circ a,x)
            =
            \delta(a,\delta(b,x))
            =
            x_2.
        \]

        \item \textbf{Use the Mealy multiplication law for outputs.}
        The output multiplication law gives
        \[
            \omega(b\circ a,x)
            =
            \omega(b,x)\circ\omega(a,\delta(b,x))
            =
            \beta\circ\gamma.
        \]

        \item \textbf{Conclude soundness.}
        Hence the derived execution judgement is interpreted by the one-step
        execution along the composite input \(b\circ a\).
    \end{enumerate}
\end{proof}

\begin{definition}[Pipeline execution rule]
\label{def:pipeline-execution-rule}
    Let
    \[
        (P,\delta_P,\omega_P):\mathbb A\rightsquigarrow\mathbb B
    \]
    and
    \[
        (Q,\delta_Q,\omega_Q):\mathbb B\rightsquigarrow\mathbb C
    \]
    be composable Mealy morphisms. Suppose
    \[
        \Gamma\vdash x:p_Px\leadsto q_Px
    \]
    and
    \[
        \Gamma\vdash y:p_Qy\leadsto q_Qy
    \]
    satisfy the compatibility condition
    \[
        q_Px=p_Qy.
    \]
    If
    \[
        \Gamma\vdash
        x\xrightarrow{\;a/\beta\;}_{P}x'
    \]
    and
    \[
        \Gamma\vdash
        y\xrightarrow{\;\beta/\gamma\;}_{Q}y',
    \]
    then the pipeline execution rule derives
    \[
        \Gamma\vdash
        (x,y)\xrightarrow{\;a/\gamma\;}_{Q\circ P}(x',y').
    \]
    Diagrammatically,
    \[
        \begin{tikzcd}[column sep=large, row sep=large]
            x
                \arrow[r, "{a/\beta}"]
            &
            x'
            \\
            y
                \arrow[r, "{\beta/\gamma}"']
            &
            y'
        \end{tikzcd}
        \qquad\leadsto\qquad
        \begin{tikzcd}[column sep=large]
            (x,y)
                \arrow[r, "{a/\gamma}"]
            &
            (x',y') .
        \end{tikzcd}
    \]
\end{definition}

\begin{proposition}[Soundness of pipeline execution]
\label{prop:soundness-pipeline-execution}
    The pipeline execution rule is sound for the composite Mealy morphism
    \[
        Q\circ P.
    \]
\end{proposition}

\begin{proof}
    \leavevmode
    \begin{enumerate}
        \item \textbf{Use the composite transition formula.}
        By Proposition~\ref{prop:pipeline-composition-mealy},
        \[
            \delta_{QP}(a,(x,y))
            =
            \bigl(
            \delta_P(a,x),
            \delta_Q(\omega_P(a,x),y)
            \bigr).
        \]

        \item \textbf{Substitute the premises.}
        Since
        \[
            x'=\delta_P(a,x),
            \qquad
            \beta=\omega_P(a,x),
        \]
        and
        \[
            y'=\delta_Q(\beta,y),
        \]
        we obtain
        \[
            \delta_{QP}(a,(x,y))=(x',y').
        \]

        \item \textbf{Use the composite output formula.}
        Again by Proposition~\ref{prop:pipeline-composition-mealy},
        \[
            \omega_{QP}(a,(x,y))
            =
            \omega_Q(\omega_P(a,x),y).
        \]
        Since
        \[
            \beta=\omega_P(a,x),
            \qquad
            \gamma=\omega_Q(\beta,y),
        \]
        the composite output is \(\gamma\).

        \item \textbf{Conclude soundness.}
        The conclusion is exactly the one-step execution judgement for
        \(Q\circ P\).
    \end{enumerate}
\end{proof}

\begin{definition}[Simulation execution rule]
\label{def:simulation-execution-rule}
    Let
    \[
        (\theta,\alpha):P\Rightarrow Q
    \]
    be a Mealy simulation, so
    \[
        \theta:Q\to P,
        \qquad
        \alpha_y:q_P(\theta y)\to q_Q(y).
    \]
    If
    \[
        \Gamma\vdash
        y\xrightarrow{\;a/\beta_Q\;}_{Q}y',
    \]
    then the simulation execution rule derives an execution
    \[
        \Gamma\vdash
        \theta y\xrightarrow{\;a/\beta_P\;}_{P}\theta y'
    \]
    together with the output comparison equation
    \[
        \alpha_y\circ \beta_P
        =
        \beta_Q\circ\alpha_{y'}.
    \]
    This is the square
    \[
        \begin{tikzcd}[column sep=large, row sep=large]
            q_P(\theta y')
                \arrow[r, "\beta_P"]
                \arrow[d, "\alpha_{y'}"']
            &
            q_P(\theta y)
                \arrow[d, "\alpha_y"]
            \\
            q_Q(y')
                \arrow[r, "\beta_Q"']
            &
            q_Q(y).
        \end{tikzcd}
    \]
\end{definition}

\begin{proposition}[Soundness of simulation execution]
\label{prop:soundness-simulation-execution}
    The simulation execution rule is sound.
\end{proposition}

\begin{proof}
    \leavevmode
    \begin{enumerate}
        \item \textbf{Use the transition part of the simulation law.}
        The transition law gives
        \[
            \theta(\delta_Q(a,y))
            =
            \delta_P(a,\theta y).
        \]
        Since
        \[
            y'=\delta_Q(a,y),
        \]
        this says
        \[
            \theta y'=\delta_P(a,\theta y).
        \]

        \item \textbf{Identify the simulated output.}
        Put
        \[
            \beta_P:=\omega_P(a,\theta y).
        \]
        Then
        \[
            \Gamma\vdash
            \theta y\xrightarrow{\;a/\beta_P\;}_{P}\theta y'.
        \]

        \item \textbf{Use the output part of the simulation law.}
        The output law gives
        \[
            \alpha_y\circ \omega_P(a,\theta y)
            =
            \omega_Q(a,y)\circ \alpha_{\delta_Q(a,y)}.
        \]
        Since
        \[
            \beta_Q=\omega_Q(a,y),
            \qquad
            y'=\delta_Q(a,y),
        \]
        this becomes
        \[
            \alpha_y\circ\beta_P
            =
            \beta_Q\circ\alpha_{y'}.
        \]
    \end{enumerate}
\end{proof}

\begin{theorem}[Soundness of the operational calculus]
\label{thm:soundness-operational-calculus}
    Let
    \[
        (P,\delta,\omega):\mathbb A\rightsquigarrow\mathbb B
    \]
    be a Mealy morphism. Each displayed judgement form is interpreted by a
    morphism in \(\mathcal E\) with the stated typing equations. Each displayed
    execution rule is sound under this interpretation. The identity,
    sequential, pipeline, and simulation rules are sound respectively by the
    Mealy unit law, the Mealy multiplication law, the pipeline-composition
    formula, and the Mealy simulation law.
\end{theorem}

\begin{proof}
    \leavevmode
    \begin{enumerate}
        \item \textbf{Interpret the basic judgements.}
        The basic judgement forms are interpreted by
        Definition~\ref{def:mealy-operational-judgements}.

        \item \textbf{Interpret one-step execution.}
        One-step execution is well typed by
        Proposition~\ref{prop:typing-one-step-execution}.

        \item \textbf{Interpret the identity rule.}
        The identity rule is sound by
        Proposition~\ref{prop:soundness-identity-execution}.

        \item \textbf{Interpret the sequential rule.}
        The sequential rule is sound by
        Proposition~\ref{prop:soundness-sequential-execution}.

        \item \textbf{Interpret the pipeline rule.}
        The pipeline rule is sound by
        Proposition~\ref{prop:soundness-pipeline-execution}.

        \item \textbf{Interpret the simulation rule.}
        The simulation rule is sound by
        Proposition~\ref{prop:soundness-simulation-execution}.

        \item \textbf{Conclude.}
        Therefore every displayed rule in the operational calculus is sound
        under the given interpretation in \(\mathcal E\).
    \end{enumerate}
\end{proof}

\begin{remark}[Conservativity]
\label{rem:operational-calculus-conservativity}
    Conversely, the one-step execution rule recovers the maps
    \[
        \delta:A_1\times_{A_0}P\to P,
        \qquad
        \omega:A_1\times_{A_0}P\to B_1.
    \]
    Thus the operational calculus is a judgemental presentation of the
    component definition of a Mealy morphism, not an additional structure.
\end{remark}

